\documentclass[12pt,a4paper]{article}

% ─── Packages ────────────────────────────────────────────────────────────────
\usepackage[T1]{fontenc}
\usepackage[utf8]{inputenc}
\usepackage{helvet}          % Arial-like font
\renewcommand{\familydefault}{\sfdefault}
\usepackage[margin=2.5cm]{geometry}
\usepackage{setspace}
\usepackage{parskip}
\usepackage{titlesec}
\usepackage{booktabs}
\usepackage{array}
\usepackage{tabularx}
\usepackage{graphicx}
\usepackage{grffile}
\usepackage{float}
\usepackage{amsmath}
\usepackage{hyperref}
\usepackage{caption}
\usepackage{fancyhdr}
\usepackage{enumitem}
\usepackage{xcolor}
\usepackage{microtype}
\usepackage{ragged2e}
\justifying

% ─── Page style ──────────────────────────────────────────────────────────────
\pagestyle{fancy}
\fancyhf{}
\rhead{CST3133 Coursework}
\lhead{CSGO Match Prediction}
\cfoot{\thepage}

% ─── Section formatting ───────────────────────────────────────────────────────
\titleformat{\section}{\large\bfseries}{}{0em}{}[\titlerule]
\titleformat{\subsection}{\normalsize\bfseries}{}{0em}{}
\titleformat{\subsubsection}{\normalsize\itshape}{}{0em}{}

\setlength{\parindent}{0pt}
\setlength{\parskip}{6pt}
\onehalfspacing

% ─── Hyperlink setup ─────────────────────────────────────────────────────────
\hypersetup{
    colorlinks=true,
    linkcolor=black,
    urlcolor=blue,
    citecolor=black
}

% ─────────────────────────────────────────────────────────────────────────────
\begin{document}

% ═══════════════════════════════════════════════════════════════════════════
%  COVER PAGE
% ═══════════════════════════════════════════════════════════════════════════
\begin{titlepage}
    \noindent
    \begin{flushright}
    
\includegraphics[height=3cm]{"Graphs/middlesexuniversity.png"}
    \end{flushright}
    \vspace{0.5cm}
    \centering
    \vspace*{0.5cm}

    {\Large \textbf{MIDDLESEX UNIVERSITY}}\\[0.3cm]
    {\normalsize The Burroughs, Hendon, London NW4 4BT}\\[2cm]

    \rule{\linewidth}{0.5pt}\\[0.5cm]
    {\LARGE \textbf{CST3133 Coursework}}\\[0.3cm]
    {\Large Applied Data Science and AI Workflow}\\[0.3cm]
    {\large Unified Report -- Dataset 1 (Structured) \& Dataset 2 (Text)}\\[0.3cm]
    \rule{\linewidth}{0.5pt}\\[2cm]

    {\large
    \begin{tabular}{ll}
        \textbf{Group Members}          & \\[0.2cm]
        Dumitru Nirca                   & M00888146 \\[0.2cm]
        Sebastian Robert Danci          & M00886707 \\[0.2cm]
        Maciej Kacper Ciba              & M00864763 \\[0.2cm]
        Peter Nagy                      & M00871555 \\[0.5cm]
        \textbf{Module}                 & CST3133 -- Advanced Topics in Data Science and AI \\[0.3cm]
        \textbf{Module Leader}          & Dr.\ Halil Yetgin \\[0.3cm]
        \textbf{Word Count}             & four thousand five hundred (Section~1) + Section~2 (max seven thousand total) \\
    \end{tabular}
    }

    \vfill
    {\normalsize Academic Year 2024/2025}
\end{titlepage}

% ═══════════════════════════════════════════════════════════════════════════
%  TABLE OF CONTENTS
% ═══════════════════════════════════════════════════════════════════════════
\newpage
\tableofcontents
\listoffigures
\listoftables
\newpage

% ═══════════════════════════════════════════════════════════════════════════
%  SECTION 1 -- DATASET 1 (STRUCTURED DATA)
% ═══════════════════════════════════════════════════════════════════════════
\section{Section 1 -- Traditional Machine Learning Workflow and Ethics (Dataset 1)}

\noindent\textit{This section covers Dataset 1: the CSGO match prediction structured dataset.}

\paragraph{Introduction and Objectives.}
This project applies a complete traditional machine learning workflow to predict CSGO match outcomes from pre-match statistics. Esports analytics increasingly informs team strategy, talent scouting, and performance benchmarking; accurate pre-match predictions can support these applications. The report documents: (one) dataset selection and problem formulation; (two) preprocessing and handling of data quality issues; (three) exploratory analysis and bias identification; (four) model development, evaluation, and interpretation; and (five) ethical considerations and mitigation strategies. We aim to achieve predictive accuracy above random chance (50\%) while maintaining interpretability and documenting limitations.

% ───────────────────────────────────────────────────────────────────────────
\subsection{1.1\quad Dataset Selection and Problem Definition (5\%)}

\subsubsection{1.1.1\quad Dataset Overview}

The \textbf{CSGO Match Prediction} dataset contains historical professional match data with comprehensive team- and player-level pre-match statistics. It is publicly available on Kaggle: \url{https://www.kaggle.com/datasets/gabrieltardochi/counter-strike-global-offensive-matches}. The dataset file \texttt{csgo\_games.csv} is submitted alongside this report. In raw form the dataset has three thousand seven hundred and eighty-seven rows and one hundred and seventy features; after preprocessing and feature engineering it retains three thousand seven hundred and sixty-three valid records and one hundred and seventy-two features for modelling.

Features are organised into two broad categories. Table~\ref{tab:features_overview} provides a detailed breakdown.

\begin{table}[H]
\centering
\caption{Exact feature inventory and descriptions.}
\label{tab:features_overview}
\begin{tabularx}{\textwidth}{lX}
\toprule
\textbf{Category} & \textbf{Features (exact names and meanings)} \\
\midrule
Team-level (raw) & \texttt{t1\_world\_rank}, \texttt{t2\_world\_rank} (world rankings 1--$\infty$); \texttt{t1\_h2h\_win\_perc}, \texttt{t2\_h2h\_win\_perc} (head-to-head win \%); \texttt{t1\_points}, \texttt{t2\_points} (removed---data leakage). \\
\midrule
Player-level (per player) & For each of \texttt{t1\_player1}--\texttt{t1\_player5} and \texttt{t2\_player1}--\texttt{t2\_player5}: \texttt{rating} (HLTV), \texttt{impact}, \texttt{kdr}, \texttt{dmr}, \texttt{kpr}, \texttt{apr}, \texttt{dpr}, \texttt{spr}, \texttt{opk\_ratio}, \texttt{opk\_rating}, \texttt{wins\_perc\_after\_fk}, \texttt{fk\_perc\_in\_wins}, \texttt{multikill\_perc}, \texttt{rating\_at\_least\_one\_perc}, \texttt{is\_sniper}, \texttt{clutch\_win\_perc}. \\
\midrule
Engineered & \texttt{t1\_avg\_rating}, \texttt{t2\_avg\_rating}; \texttt{t1\_avg\_impact}, \texttt{t2\_avg\_impact}; \texttt{t1\_avg\_kdr}, \texttt{t2\_avg\_kdr}; \texttt{rank\_diff}; \texttt{h2h\_advantage}. \\
\bottomrule
\end{tabularx}
\end{table}

\textbf{Target variable:} \texttt{winner}---binary (Team~1 or Team~2). Class distribution is near-balanced (51.05\% / 48.95\%), so no oversampling or class weighting is required.

\subsubsection{1.1.2\quad Problem Definition}

This project addresses a \textbf{binary classification problem}: predicting which team will win a CSGO match given only pre-match available statistics.

\textbf{Specific use cases:} (one) Team management and scouting---front offices can use predictions to prioritise opponent analysis and benchmark expected win probability against actual results; (two) tournament organisers can inform seeding decisions and scheduling; (three) content and broadcast partners can generate pre-match storylines and tailor commentary; (four) sponsorship and investment teams can assess team strength quantitatively. The problem has a \textbf{gambling-like nature}: like football matches or roulette, outcomes are inherently stochastic and not fully predictable from available inputs. The model provides probability estimates, not deterministic forecasts; interpretability and documented limitations are essential for responsible deployment.

% ───────────────────────────────────────────────────────────────────────────
\subsection{1.2\quad Data Preprocessing (10\%)}

\subsubsection{1.2.1\quad Handling Missing Values and Noise}

The dataset contained naturally occurring missing values (six thousand seven hundred and forty-three NaNs across thirty-four player-level columns); additional NaN values were introduced in selected columns per coursework requirements to create a noisy version for preprocessing. The focus is on identifying and mitigating these data quality issues. \textbf{Imputation strategy:} Median imputation for numerical columns, as the median is robust to outliers in player statistics where exceptional performances can skew means; mode imputation for categorical columns. Figure~\ref{fig:missing} shows missing values by column before imputation.

Table~\ref{tab:before_after} summarises before and after preprocessing.

\begin{table}[H]
\centering
\caption{Data quality: before and after preprocessing.}
\label{tab:before_after}
\begin{tabularx}{\textwidth}{l>{\raggedright\arraybackslash}X>{\centering\arraybackslash}X}
\toprule
\textbf{Aspect} & \textbf{Before (with NaN / noise)} & \textbf{After (processed)} \\
\midrule
Rows & 3,787 & 3,763 \\
Columns / features & 170 & 172 (numerical only) \\
Missing values & 6,743+ (natural + introduced; imputed) & 0 \\
Outliers & Present (median imputation) & Retained, scaled \\
Leakage features & \texttt{t1\_points}, \texttt{t2\_points}, \texttt{points\_diff} & Removed \\
Train / test split & -- & 3,010 / 753 \\
\bottomrule
\end{tabularx}
\end{table}

Outliers were identified via boxplot inspection of key features (e.g., \texttt{rank\_diff}, ratings), right-skewed distributions in world rankings, and domain knowledge---exceptional player performances (e.g., KDR above two) are valid but rare. Median imputation was chosen for robustness to extremes. Range validation confirmed no impossible values (e.g., negative percentages). After imputation, the dataset contains zero missing values.

\begin{figure}[H]
\centering
\includegraphics[width=0.9\textwidth]{"Graphs/missing_values_by_column.png"}
\caption{Missing values by column before imputation; guided imputation strategy.}
\label{fig:missing}
\end{figure}

The \texttt{csgo\_games.csv} file and preprocessing documentation are submitted alongside this report.

\subsubsection{1.2.2\quad Feature Engineering and Scaling}

\textbf{Eight engineered features:}

\begin{table}[H]
\centering
\caption{Engineered features and their rationale.}
\begin{tabularx}{\textwidth}{lX}
\toprule
\textbf{Feature} & \textbf{Description and Rationale} \\
\midrule
\texttt{t1\_avg\_rating}, \texttt{t2\_avg\_rating} & Average HLTV rating across each team's five players. Player rating is the primary individual performance metric in professional CSGO. \\
\texttt{t1\_avg\_impact}, \texttt{t2\_avg\_impact} & Average impact score per team, measuring round-winning contributions beyond kills. \\
\texttt{t1\_avg\_kdr}, \texttt{t2\_avg\_kdr} & Average Kill-Death Ratio per team, a standard measure of individual survivability and lethality. \\
\texttt{rank\_diff} & Difference in world rankings (\texttt{t1\_world\_rank} $-$ \texttt{t2\_world\_rank}). Captures the relative competitive strength gap. \\
\texttt{h2h\_advantage} & Difference in head-to-head win percentages. Captures historical team-versus-team performance patterns. \\
\bottomrule
\end{tabularx}
\end{table}

Three features (\texttt{t1\_points}, \texttt{t2\_points}, \texttt{points\_diff}) were removed to prevent \textbf{data leakage}, as tournament points are updated following match results and would not be available prior to a match. All one hundred and seventy-two retained numerical features were standardised using \textbf{StandardScaler}; the target was encoded with \textbf{LabelEncoder}. The final split is three thousand and ten training samples and seven hundred and fifty-three test samples (stratified eighty/twenty).

% ───────────────────────────────────────────────────────────────────────────
\subsection{1.3\quad Exploratory Data Analysis (10\%)}

\subsubsection{1.3.1\quad Class Distribution}

As shown in the dataset overview, the target variable is near-perfectly balanced (51.05\% / 48.95\%). This ensures that standard accuracy is a valid primary evaluation metric and that models are not inadvertently rewarded for predicting the majority class.

\subsubsection{1.3.2\quad Feature Distributions}

Figure~\ref{fig:class_dist} shows the class distribution. Histograms of key features (Figures~\ref{fig:dist_key_features}) reveal: world rankings exhibit a right-skewed distribution reflecting concentration among top-ranked teams; average ratings cluster near one (the HLTV baseline); \texttt{rank\_diff} spans a wide range, indicating diverse competitive matchups.

\begin{figure}[H]
\centering
\includegraphics[width=0.75\textwidth]{"Graphs/Distribution of Match Winners.png"}
\caption{Match outcome distribution (t1 vs.\ t2); nearly balanced.}
\label{fig:class_dist}
\end{figure}


\begin{figure}[H]
\centering
\includegraphics[width=0.95\textwidth]{"Graphs/Distribution_of_key_features.png"}
\caption{Key feature histograms: rankings, ratings, \texttt{rank\_diff}, \texttt{h2h\_advantage}.}
\label{fig:dist_key_features}
\end{figure}

\subsubsection{1.3.3\quad Feature Dependencies}

Correlation analysis (Figure~\ref{fig:correlation}) reveals important feature dependencies: \texttt{t1\_avg\_rating} and \texttt{t2\_avg\_rating} correlate strongly with individual player ratings (expected, since they are team averages); \texttt{rank\_diff} correlates with team-average ratings. This multicollinearity motivates feature selection to reduce redundant predictors.

\begin{figure}[H]
\centering
\includegraphics[width=0.9\textwidth]{"Graphs/Corellation Features.png"}
\caption{Correlation heatmap; informed feature selection.}
\label{fig:correlation}
\end{figure}

\subsubsection{1.3.4\quad Key Trends}

When Team~1 holds a better world ranking (more negative \texttt{rank\_diff}), Team~1 wins more frequently---validating world ranking as a dominant predictive signal. \texttt{rank\_diff} and team-average ratings dominate; player-level statistics (clutch percentage, multi-kill percentage) contribute. Figures~\ref{fig:rank_diff_winner} and \ref{fig:team_ratings} visualise rank difference and team ratings by match outcome.

\begin{figure}[H]
\centering
\includegraphics[width=0.95\textwidth]{"Graphs/Rank_difference_by_winner_and_rank_distribtuion_by_winner.png"}
\caption{\texttt{rank\_diff} by winner; supports it as strong predictor.}
\label{fig:rank_diff_winner}
\end{figure}


\begin{figure}[H]
\centering
\includegraphics[width=0.85\textwidth]{"Graphs/team 1 Rating_ Team 2 Rating.png"}
\caption{Team ratings by outcome; higher ratings correlate with wins.}
\label{fig:team_ratings}
\end{figure}

\subsubsection{1.3.5\quad Potential Biases Identified in the Data}

\begin{enumerate}[leftmargin=1.5em]
    \item \textbf{Temporal Bias:} The dataset spans multiple competitive seasons. CSGO meta-game shifts (strategy, weapon balance, map pool) mean that patterns from earlier years may not generalise to more recent matches.
    \item \textbf{Team Representation Bias:} Top-ranked teams compete more frequently and therefore dominate the dataset. The model may be poorly calibrated for underrepresented lower-ranked or regional teams.
    \item \textbf{Ranking System Bias:} World rankings are a lagging indicator and do not immediately reflect roster changes or strategic pivots, introducing systematic noise in the most predictive feature.
\end{enumerate}

% ───────────────────────────────────────────────────────────────────────────
\subsection{1.4\quad Model Development and Evaluation (15\%)}

\subsubsection{1.4.1\quad Models Trained}

Three model categories were developed:

\paragraph{Logistic Regression.} A linear baseline classifier modelling the log-odds as $\log\frac{P(y=1|\mathbf{x})}{1-P(y=1|\mathbf{x})} = \mathbf{w}^\top \mathbf{x} + b$. It is interpretable, computationally efficient, and provides calibrated probability estimates---the natural baseline for this problem.

\paragraph{Decision Tree.} A non-linear classifier using Gini impurity $G(t) = 1 - \sum_c p(c|t)^2$ to recursively partition the feature space. Trained with \texttt{max\_depth}=10 and \texttt{min\_samples\_split}=20 to limit overfitting. Produces direct feature importance scores based on impurity reduction.

\paragraph{K-Means.} Per the coursework requirement for an unsupervised method, K-Means with $k=2$ was applied to test whether match outcomes can be recovered from natural groupings in the feature space. The objective minimises $\sum_i \sum_{\mathbf{x} \in C_i} \|\mathbf{x} - \boldsymbol{\mu}_i\|^2$. The result (ARI $\approx 0$) demonstrates that unsupervised clustering cannot separate winners, justifying supervised classification.

\subsubsection{1.4.2\quad Evaluation Metrics}

For classification models:

\begin{itemize}[leftmargin=1.5em]
    \item \textbf{Accuracy:} $\frac{TP + TN}{TP + TN + FP + FN}$ --- proportion of correctly classified matches.
    \item \textbf{Precision:} $\frac{TP}{TP + FP}$ — of all predicted Team~1 wins, how many were correct.
    \item \textbf{Recall:} $\frac{TP}{TP + FN}$ — of all actual Team~1 wins, how many were identified.
    \item \textbf{F1-Score:} $\frac{2 \cdot \text{Precision} \cdot \text{Recall}}{\text{Precision} + \text{Recall}}$ — harmonic mean balancing precision and recall.
\end{itemize}

For K-Means:
\begin{itemize}[leftmargin=1.5em]
    \item \textbf{Adjusted Rand Index (ARI):} $\text{ARI} = \frac{\text{RI} - \mathbb{E}[\text{RI}]}{\max(\text{RI}) - \mathbb{E}[\text{RI}]}$, where RI is the Rand Index. Agreement between cluster assignments and true labels, corrected for chance. ARI $\approx 0$ indicates random assignment.
    \item \textbf{Silhouette Score:} $s(i) = \frac{b(i) - a(i)}{\max(a(i), b(i))}$, where $a(i)$ is mean intra-cluster distance and $b(i)$ is mean nearest-cluster distance. Values near 1 indicate well-separated clusters.
\end{itemize}

\subsubsection{1.4.3\quad Results}

Table~\ref{tab:results} presents the performance of all trained models on the held-out test set.

\begin{table}[H]
\centering
\caption{Model performance summary on the test set.}
\label{tab:results}
\begin{tabular}{lcccc}
\toprule
\textbf{Model} & \textbf{Accuracy} & \textbf{Precision} & \textbf{Recall} & \textbf{F1-Score} \\
\midrule
Logistic Regression           & 55.91\%  & 0.5476 & 0.5772 & 0.5620 \\
Decision Tree                 & 53.78\%  & 0.5258 & 0.5799 & 0.5515 \\
K-Means (ARI / Sil.)          & $-$0.0003 & -- & -- & -- / 0.0645 \\
LR with Feature Selection (best) & 58.70\%  & 0.5751 & 0.6016 & 0.5881 \\
\bottomrule
\end{tabular}
\end{table}

\paragraph{Confusion matrix.}
Table~\ref{tab:cm_lr}: baseline LR; rows = true, columns = predicted.

\begin{table}[H]
\centering
\caption{Confusion matrix: Logistic Regression (baseline) on test set.}
\label{tab:cm_lr}
\begin{tabular}{cc|cc}
 & & \multicolumn{2}{c}{\textbf{Predicted}} \\
 & & Team~2 (0) & Team~1 (1) \\
\hline
\textbf{True} & Team~2 (0) & TN = 221 & FP = 163 \\
 & Team~1 (1) & FN = 145 & TP = 224 \\
\end{tabular}
\end{table}


\begin{figure}[H]
\centering
\includegraphics[width=0.6\textwidth]{"Graphs/Confuxion_Matrix_Logistic_Reversion.png"}
\caption{Confusion matrix: Logistic Regression baseline.}
\label{fig:cm_lr}
\end{figure}

Models are \textbf{marginally better than random} (50\%): baseline LR achieves 55.91\% and the best model 58.70\%, an improvement of roughly nine percentage points. Near-equal false positive and false negative counts indicate weak class separation. The Decision Tree (Figure~\ref{fig:cm_dt}) performs similarly but with slightly lower accuracy.

\begin{figure}[H]
\centering
\includegraphics[width=0.6\textwidth]{"Graphs/Confusion_Matrix_Decision_Tree.png"}
\caption{Confusion matrix: Decision Tree.}
\label{fig:cm_dt}
\end{figure}


\begin{figure}[H]
\centering
\includegraphics[width=0.95\textwidth]{"Graphs/All Models Accuracy Comparison _ Top 5 Models.png"}
\caption{Top five models by accuracy; LR+FS best (58.70\%).}
\label{fig:all_models}
\end{figure}

\begin{figure}[H]
\centering
\includegraphics[width=0.95\textwidth]{"Graphs/All Models Accuracy Comparison _ Top 5 Models - Metrics.png"}
\caption{Top five models: precision, recall, F1.}
\label{fig:all_models_metrics}
\end{figure}


\begin{figure}[H]
\centering
\includegraphics[width=0.95\textwidth]{"Graphs/Model Accuracy Comparison Top 3 Models.png"}
\caption{Top three models comparison.}
\label{fig:top3_models}
\end{figure}

\begin{figure}[H]
\centering
\includegraphics[width=0.95\textwidth]{"Graphs/Model Accuracy Comparison and Detailed Metrics Comparison.png"}
\caption{All models: accuracy and metrics.}
\label{fig:model_metrics}
\end{figure}


\begin{figure}[H]
\centering
\includegraphics[width=0.95\textwidth]{"Graphs/Optimized Models Accuracy and Top 3 Optimized.png"}
\caption{Optimised models post-tuning.}
\label{fig:optimized_models}
\end{figure}


\begin{figure}[H]
\centering
\includegraphics[width=0.9\textwidth]{"Graphs/Actual Classes(PCA Visualization) K-Means Clusters (PCA Visualization).png"}
\caption{PCA: true outcomes vs.\ K-Means clusters ($k=2$); ARI $\approx 0$.}
\label{fig:pca_kmeans}
\end{figure}

\paragraph{Interpretation.}
The best-performing model, \textbf{LR with Feature Selection} at 58.70\%, improves over the baseline by 2.79\%. This must be interpreted \textbf{critically}: 58.7\% accuracy is barely above random chance (50\% for balanced binary classification). The modest gain reflects the fundamental unpredictability of match outcomes---a problem with a \textbf{gambling-like nature}, akin to football matches or roulette, where outcomes are inherently stochastic. Esports literature typically reports pre-match accuracies in the 55--65\% range; our results sit within this band. Despite trying over sixty model configurations, ensembles, and hyperparameter combinations, we could not exceed approximately 60\% on this dataset; further gains would likely require a different dataset (e.g., in-game state, map-specific data). K-Means (ARI $\approx 0$) confirms that unsupervised methods cannot recover outcomes.

\subsubsection{1.4.4\quad Model Improvements}

Substantial effort was devoted to hyperparameter tuning via \texttt{GridSearchCV} and \texttt{RandomizedSearchCV} with three-fold cross-validation. Table~\ref{tab:hyperparameter_tuning} documents the parameter grids, search methods, and best configurations.

\begin{table}[H]
\centering
\caption{Systematic hyperparameter tuning: parameter grids and best configurations.}
\label{tab:hyperparameter_tuning}
\small
\renewcommand{\arraystretch}{1.6}
\begin{tabular}{>{\raggedright\arraybackslash}p{3.4cm}>{\raggedright\arraybackslash}p{5.2cm}>{\raggedright\arraybackslash}p{4.2cm}>{\raggedleft\arraybackslash}p{2.6cm}}
\toprule
\textbf{Model} & \textbf{Parameter grid} & \textbf{Best parameters} & \textbf{CV / test} \\
\midrule
\mbox{Logistic Regression} & $C \in \{0.01, 0.1, 1, 10\}$; \texttt{class\_weight} & $C = 0.01$ (stronger regularisation) & 58.30\% \\
\mbox{Random Forest} & \texttt{n\_estimators} $\in \{50, 100\}$, \texttt{max\_depth} $\in \{10, 15, 20\}$, \texttt{min\_samples\_split} $\in \{10, 20\}$ & \texttt{max\_depth}=10, \texttt{min\_samples\_split}=10, \texttt{n\_estimators}=100 & 58.57 / 58.03\% \\
\mbox{XGBoost} & \texttt{n\_estimators} $\in \{100, 200\}$, \texttt{max\_depth} $\in \{4, 6, 8\}$, \texttt{learning\_rate} $\in \{0.05, 0.1\}$, \texttt{subsample} $\in \{0.8, 0.9\}$ & \texttt{learning\_rate}=0.05, \texttt{max\_depth}=4, \texttt{n\_estimators}=100, \texttt{subsample}=0.8 & 58.07\% \\
\mbox{LightGBM} & \texttt{n\_estimators} $\in \{200, 300\}$, \texttt{max\_depth} $\in \{5, 7, 9\}$, \texttt{learning\_rate} $\in \{0.03, 0.05, 0.1\}$, \texttt{subsample} $\in \{0.8, 0.9\}$ & \texttt{learning\_rate}=0.03, \texttt{max\_depth}=5, \texttt{n\_estimators}=200, \texttt{subsample}=0.8 & 56.98\% \\
\midrule
\multicolumn{4}{>{\raggedright\arraybackslash}p{15.2cm}}{\textit{Extended tuning:} \texttt{RandomizedSearchCV} for XGBoost and LightGBM with expanded grids (\texttt{colsample\_bytree}, \texttt{num\_leaves}, etc.). Decision threshold sweep over 0.38--0.63 on out-of-fold probabilities.} \\
\multicolumn{4}{>{\raggedright\arraybackslash}p{15.2cm}}{\textit{Note:} Candidates evaluated---RF: 12, XGBoost: 24, LightGBM: 36 (three-fold CV).} \\
\bottomrule
\end{tabular}
\end{table}

Despite extensive tuning, \textbf{LR with Feature Selection} (58.70\%) remained the best model; tuned Random Forest (58.03\%) came closest. Techniques applied: \texttt{SelectFromModel} based on Random Forest importance (retaining the top fifty features); LR with Feature Selection; tuned Random Forest, XGBoost, and LightGBM via GridSearchCV and RandomizedSearchCV; Soft Voting and Stacking ensembles; decision threshold optimisation over 0.38--0.63. None exceeded LR with Feature Selection.

\subsubsection{1.4.5\quad Feature Importance}

Table~\ref{tab:features} presents the top 5 most predictive features identified by the Decision Tree model.

\begin{table}[H]
\centering
\caption{Top 5 most important features from Decision Tree.}
\label{tab:features}
\begin{tabular}{lll}
\toprule
\textbf{Rank} & \textbf{Feature} & \textbf{Importance} \\
\midrule
1 & \texttt{rank\_diff}                    & 0.0836 \\
2 & \texttt{t1\_avg\_rating}               & 0.0280 \\
3 & \texttt{t1\_player2\_clutch\_win\_perc} & 0.0267 \\
4 & \texttt{t2\_avg\_rating}               & 0.0262 \\
5 & \texttt{t2\_player5\_clutch\_win\_perc} & 0.0260 \\
\bottomrule
\end{tabular}
\end{table}

\begin{figure}[H]
\centering
\includegraphics[width=0.9\textwidth,height=0.32\textheight,keepaspectratio]{"Graphs/Top_15_Feature_Importance.png"}
\caption{Top fifteen features: Decision Tree importance.}
\label{fig:top15_features}
\vspace{0.3em}
\includegraphics[width=0.9\textwidth,height=0.32\textheight,keepaspectratio]{"Graphs/Top 20 Feature Importance - Basic Catboost.png"}
\caption{Top twenty features: CatBoost; aligns with Decision Tree.}
\label{fig:top20_catboost}
\end{figure}

\texttt{rank\_diff} and team-average ratings dominate; Figures~\ref{fig:top15_features} and \ref{fig:top20_catboost} confirm cross-model agreement.

\subsubsection{1.4.6\quad Interpretation}

Logistic Regression outperforming the Decision Tree indicates that the relationship between features and match outcomes is largely linear. Over sixty configurations were tried across ten model types (Logistic Regression, Decision Tree, Random Forest, Gradient Boosting, XGBoost, LightGBM, CatBoost, Voting Classifier, Stacking Classifier, and K-Means); the best remained LR with Feature Selection at 58.70\%. \textbf{Critical assessment:} The problem has a \textbf{gambling-like nature}---like predicting football match results or roulette spins, outcomes are inherently stochastic. Predictions are indicative rather than definitive; the model must \textbf{not} be deployed for betting or gambling purposes (Section~1.5).

% ───────────────────────────────────────────────────────────────────────────
\subsection{1.5\quad Ethical Considerations (5\%)}

\subsubsection{1.5.1\quad Potential Biases}

\textbf{1. Temporal:} The dataset spans multiple years; CSGO meta evolves through patches and map pool changes---a concept drift risk. \textbf{2. Team representation:} Top-ranked teams compete more frequently and dominate the dataset; the model may underperform on lower-ranked or regional teams. \textbf{3. Ranking:} \texttt{rank\_diff} relies on world rankings, which are lagging indicators; roster or strategy changes may not be reflected. \textbf{4. Data collection:} Major international tournaments are over-represented; regional leagues underrepresented. \textbf{5. Aggregation:} Team-level averages may miss team chemistry and role specialisation.

\subsubsection{1.5.2\quad Mitigation Strategies}

\textbf{1. Temporal:} Implement time-based cross-validation; retrain with post-patch data; add recency features (e.g., rolling rating over the last ten matches). \textbf{2. Representation:} Apply stratified sampling by team tier and region when constructing train/test sets. \textbf{3. Rankings:} Prioritise relative features (\allowbreak\texttt{rank\_diff}, rating difference) rather than absolute values. \textbf{4. Data:} Expand to regional leagues; document which tournaments and date ranges are covered. \textbf{5. Auditing:} Maintain interpretable models; conduct fairness audits comparing accuracy across tiers and regions.

\textbf{6. Gambling:} The model must \textbf{not} be used for betting---the problem is inherently stochastic (like football or roulette); 58.7\% versus 50\% is marginal. Communicate uncertainty; comply with GDPR for player data.

\textbf{7. Transparency:} Document HLTV sources, preprocessing steps, and limitations; share with stakeholders (analysts, broadcasters, organisers) so they can judge when to trust predictions versus rely on expert judgment.

\subsubsection{1.5.3\quad Key Recommendations}

Key recommendations: (one) use the best-performing model, LR with Feature Selection, for match predictions; (two) retrain with recent data to account for meta-game changes; (three) monitor performance across different teams and time periods; (four) validate predictions with domain experts before deployment. Do not deploy for betting or gambling purposes.

% ═══════════════════════════════════════════════════════════════════════════
%  SECTION 2 -- DATASET 2 (TEXT / NLP / DEEP LEARNING)
% ═══════════════════════════════════════════════════════════════════════════
\newpage
\section{Section 2 -- Natural Language Processing and Deep Learning (Dataset 2)}

\noindent\textit{This section covers Dataset 2: sentiment analysis of CSGO Steam reviews to classify reviews as positive or negative. The dataset is publicly available on Kaggle: \url{https://www.kaggle.com/datasets/noahx1/csgo-steam-reviews}.}

\subsection{2.1\quad Text Dataset Selection and Preprocessing (15\%)}

The CSGO Steam Reviews dataset will be used for sentiment analysis: classifying reviews as positive or negative to derive an overall rating or sentiment score. Preprocessing will include: text cleaning, tokenisation, stopword removal, and pre-trained embeddings (e.g., GloVe, Word2Vec). Dataset source: \url{https://www.kaggle.com/datasets/noahx1/csgo-steam-reviews}.

\subsection{2.2\quad Deep Learning Model Implementation (20\%)}

\noindent\textit{[To be completed: Neural network design (RNN, LSTM); layer descriptions---embeddings, recurrent, fully connected, dropout, early stopping; hyperparameter tuning.]}

\subsection{2.3\quad Evaluation and Insights (10\%)}

\noindent\textit{[To be completed: Evaluation metrics---confusion matrix, accuracy, precision, recall; loss curves; strengths, limitations, and areas for improvement.]}

% ═══════════════════════════════════════════════════════════════════════════
%  REFERENCES (Harvard style, alphabetical order)
% ═══════════════════════════════════════════════════════════════════════════
\newpage
\section{References}

\begin{list}{}{%
  \setlength{\leftmargin}{2em}%
  \setlength{\itemindent}{-2em}%
  \setlength{\itemsep}{0.75em}%
  \setlength{\parsep}{0pt}}

\item Chen, T. and Guestrin, C. (2016) `XGBoost: A Scalable Tree Boosting System', \textit{Proceedings of the 22nd ACM SIGKDD International Conference on Knowledge Discovery and Data Mining}, pp.\ 785--794. Available at: \url{https://doi.org/10.1145/2939672.2939785} (Accessed: 7 February 2025).

\item Harris, C.R., Millman, K.J., van der Walt, S.J. \textit{et al}. (2020) `Array programming with NumPy', \textit{Nature}, 585, pp.\ 357--362. Available at: \url{https://doi.org/10.1038/s41586-020-2649-2} (Accessed: 7 February 2025).

\item HLTV.org (2024) \textit{CSGO Match Statistics}. Available at: \url{https://www.hltv.org} (Accessed: 7 February 2025).

\item Kaggle (2024) \textit{Counter-Strike: Global Offensive matches}. Available at: \url{https://www.kaggle.com/datasets/gabrieltardochi/counter-strike-global-offensive-matches} (Accessed: 7 February 2025).

\item Kaggle (2024) \textit{CSGO Steam Reviews}. Available at: \url{https://www.kaggle.com/datasets/noahx1/csgo-steam-reviews} (Accessed: 7 February 2025).

\item McKinney, W. (2010) `Data structures for statistical computing in Python', in van der Walt, S. and Millman, J. (eds) \textit{Proceedings of the 9th Python in Science Conference}, pp.\ 56--61.

\item Pedregosa, F., Varoquaux, G., Gramfort, A. \textit{et al}. (2011) `Scikit-learn: Machine learning in Python', \textit{Journal of Machine Learning Research}, 12, pp.\ 2825--2830. Available at: \url{https://jmlr.org/papers/v12/pedregosa11a.html} (Accessed: 7 February 2025).

\end{list}

% ═══════════════════════════════════════════════════════════════════════════
\end{document}
